\section{Nhiệm vụ 2. Phương pháp mới và biến đổi đồ thị}
\subsection{Task A. Dự đoán quan hệ trên GQA}

Task A áp dụng lại bài toán PredCls của Nhiệm vụ 1 trên dataset thứ hai,
GQA~\cite{hudson2019gqa}. Với bounding box và nhãn đối tượng đã biết, mô hình
vẫn dự đoán predicate cho từng cặp có hướng. Trên GQA, nhóm khảo sát liệu
việc biểu diễn mỗi predicate bằng một vector đại diện (\textit{prototype})
có giúp nhận diện quan hệ trên nhiều lớp hơn hay chủ yếu cải thiện recall
trên toàn ảnh.

\subsubsection{Dataset và split}
Nhóm chia phần train chính thức theo ảnh thành 69.942 ảnh train và
5.000 ảnh validation, dùng 10.696 ảnh từ phần validation chính thức làm
test của đồ án. Vocabulary xây từ annotation của riêng tập train có
310 foreground predicates. Hình~\ref{fig:gqa-train-predicates} cho thấy
hai quan hệ trái/phải chiếm 89,16\% số chú thích train, trong khi 42 lớp
có dưới 20 chú thích mỗi lớp. Phân bố tương tự xuất hiện trên test. Hai nhãn
\textit{to the left of} và \textit{to the right of} chiếm tổng cộng
476.716 trong 534.437 chú thích quan hệ (89,2\%). R@$K$ đo tỷ lệ
quan hệ chuẩn được tìm thấy trong top-$K$ của từng ảnh rồi lấy trung bình
trên 10.298 ảnh có quan hệ chuẩn. Vì vậy, trên ảnh có nhiều chú thích
trái/phải, tìm đúng hai quan hệ này có thể đem lại R@$K$ cao dù mô hình
bỏ sót các predicate khác. Với mR@$K$, nhóm tính recall theo ảnh cho từng
predicate rồi lấy trung bình đều trên 310 lớp. Hai lớp trái/phải chỉ chiếm
$2/310$ trọng số của mR@$K$. Có 15 lớp vắng trong test được tính recall bằng 0.

\begin{figure}[!ht]\centering
\begin{tikzpicture}
\begin{semilogyaxis}[report plot,width=.98\linewidth,height=5.4cm,
 xmin=0,xmax=315,ymin=1,ymax=4000000,xtick={1,100,200,310},
 ytick={1,100,10000,1000000},ymajorgrids,
 xlabel={Hạng quan hệ trên tập train},ylabel={Số chú thích (log)}]
\path[fill=mutedblue] (axis cs:268.5,1) rectangle (axis cs:310.5,4000000);
\draw[dashed,cred] (axis cs:1,20)--(axis cs:310,20);
\addplot[inkblue,thin,mark=*,mark size=.42pt] table[x=rank,y=count]{figures/pgf/gqa-frequency.tsv};
\draw[inkblue] (axis cs:2,1578919)--(axis cs:20,1150000)
 node[right,font=\scriptsize] {\textit{to the left of / to the right of}};
\draw[inkblue] (axis cs:3,84299)--(axis cs:32,55000) node[right] {\textit{on}};
\node[above right,cred] at (axis cs:20,20) {20 chú thích};
\draw[inkblue] (axis cs:309,2)--(axis cs:237,2.6) node[left] {\textit{higher than} (2)};
\end{semilogyaxis}
\end{tikzpicture}

\caption{Phân bố 3.541.725 chú thích quan hệ trên GQA train. Mỗi điểm
ứng với một trong 310 predicate, xếp theo tần suất giảm dần. Trục tung dùng
thang log. Hai lớp đầu có 1.578.919 chú thích mỗi lớp. Vùng tô ở đuôi
gồm 42 lớp có dưới 20 chú thích.}
\label{fig:gqa-train-predicates}
\end{figure}

\subsubsection{Baseline và điều chỉnh đầu vào của PE-Net}
Nhóm dùng MLP trên cặp đối tượng làm baseline phi đồ thị, GAT để thêm
message passing và GAT + KG để khảo sát nhúng WordNet--TransE.
Để đáp ứng yêu cầu triển khai một phương pháp từ năm 2023, nhóm chọn
PE-Net (CVPR 2023)~\cite{zheng2023penet}. Phương pháp học biểu diễn cặp
đối tượng và prototype của predicate trong cùng một không gian.
Nhóm điều chỉnh đầu vào để dùng đặc trưng CLIP ViT-B/32 đã trích sẵn.
Mỗi đối tượng được biểu diễn bằng 512 giá trị CLIP của vùng ảnh và năm giá trị
hình học chuẩn hóa (tọa độ tâm, chiều rộng, chiều cao và diện tích của box).
PE-Net gốc lấy đặc trưng từ vùng bao cả hai đối tượng (\textit{union box}).
Trong triển khai của nhóm, một MLP nhận hai vector 517 chiều ghép theo thứ tự
subject rồi object và tạo vector 256 chiều cho cặp. Mô hình khởi tạo
embedding tên lớp bằng GloVe và chấm điểm từng predicate bằng cosine
similarity giữa biểu diễn cặp và prototype. PE-Net gốc báo cáo trên Visual
Genome và Open Images. Các điểm dưới đây dùng đặc trưng, vocabulary và
split riêng của GQA trong đồ án.

Ba cấu hình dùng cùng cách biểu diễn để tách tác dụng của loss.
\textit{PE-CE} học bằng cross-entropy. Khi một cặp có nhiều predicate
chuẩn, xác suất mục tiêu chia đều cho các nhãn đó.
\textit{PE-PL} thêm thành phần margin của Prototype-guided Learning,
khuyến khích biểu diễn cặp gần prototype đúng hơn các prototype âm gần nhất.
\textit{PE-PLPR} thêm Prototype Regularization để tách các prototype
của các lớp khác nhau. Đường trên của Hình~\ref{fig:gqa-penet-aligng}
minh họa cách PE-Net chấm điểm cặp \textit{person--jeans}.

\begin{figure*}[!t]\centering
\includegraphics[width=.98\textwidth]{gqa-penet-aligng.pdf}
\caption{PE-Net và nhánh AlignG trên ảnh GQA validation 1056. Đường trên
mã hóa cặp có hướng và so khớp với prototype toàn cục. Khi bật AlignG,
các cặp trong cùng ảnh góp thông tin để cập nhật prototype bằng GRU.
Attention ngược dùng các prototype theo ảnh để hiệu chỉnh vector cặp
trước phép so khớp cosine. PE-PLPR và PE-PLPR + AlignG (seed 42) đều cho Pred là \textit{wearing} ở hạng 7, khớp với GT. \gtpredlegend}
\label{fig:gqa-penet-aligng}
\end{figure*}

\subsubsection{Kết quả test trên GQA}
Bảng~\ref{tab:gqatest} tổng hợp đủ 27 lượt của chín cấu hình. Mỗi lượt dùng
checkpoint đã chọn theo GC mR@50 trên validation.
PE-PL đạt R@50 và R@100 cao nhất (78,21\% và 88,48\%). Khi thêm PR vào
PE-PL, mR@50 tăng từ 3,40 lên 4,07\%, còn R@50 giảm từ 78,21 xuống
78,00\%. Với baseline, chuyển từ MLP sang GAT làm R@50 tăng từ 76,86 lên
77,46\% nhưng mR@50 giảm từ 3,58 xuống 2,97\%.
Phân tích theo predicate cho thấy ảnh hưởng của phân bố nhãn. Với PE-PL,
recall ở top-50 trung bình trên hai lớp trái/phải đạt 86,51\%, còn trung bình
trên 308 lớp còn lại chỉ đạt 2,86\% (trung bình ba seed, giữ các lớp vắng
trong test theo cùng quy ước).

\begin{table*}[!t]\centering\small
\caption{PredCls trên GQA test của đồ án (10.696 ảnh). Điểm tính theo \%,
trung bình $\pm$ độ lệch chuẩn mẫu qua ba seed 42, 123 và 3407.}
\label{tab:gqatest}
\begin{tabular}{@{}lrrrr@{}}\toprule
Cấu hình & R@50 & R@100 & mR@50 & mR@100\\\midrule
MLP & $76{,}86\pm0{,}45$ & $87{,}39\pm0{,}29$ & $3{,}58\pm0{,}38$ & $5{,}88\pm0{,}46$ \\
GAT (CE) & $77{,}46\pm0{,}11$ & $87{,}86\pm0{,}10$ & $2{,}97\pm0{,}24$ & $4{,}88\pm0{,}19$ \\
GAT (adjusted loss) & $77{,}07\pm0{,}05$ & $87{,}48\pm0{,}04$ & $\mathbf{4{,}62\pm0{,}71}$ & $\mathbf{7{,}97\pm0{,}86}$ \\
GAT + KG & $77{,}42\pm0{,}04$ & $87{,}84\pm0{,}05$ & $3{,}04\pm0{,}19$ & $4{,}85\pm0{,}22$ \\
PE-CE & $78{,}11\pm0{,}17$ & $88{,}40\pm0{,}13$ & $2{,}93\pm0{,}16$ & $4{,}68\pm0{,}39$ \\
PE-PL & $\mathbf{78{,}21\pm0{,}09}$ & $\mathbf{88{,}48\pm0{,}03}$ & $3{,}40\pm0{,}22$ & $5{,}48\pm0{,}37$ \\
PE-PLPR & $78{,}00\pm0{,}03$ & $88{,}24\pm0{,}05$ & $4{,}07\pm0{,}20$ & $6{,}78\pm0{,}18$ \\
PE-CE + AlignG & $78{,}03\pm0{,}36$ & $88{,}26\pm0{,}35$ & $3{,}52\pm0{,}32$ & $5{,}68\pm0{,}45$ \\
Edge-only & $77{,}79\pm0{,}02$ & $88{,}00\pm0{,}02$ & $4{,}31\pm0{,}19$ & $7{,}00\pm0{,}17$ \\
PE-PLPR + AlignG & $78{,}06\pm0{,}07$ & $88{,}22\pm0{,}06$ & $\mathbf{4{,}62\pm0{,}57}$ & $7{,}44\pm0{,}81$ \\
\bottomrule

\end{tabular}\\[2pt]
{\footnotesize GC = \textit{Graph Constraint}. Mỗi cặp có hướng giữ một
predicate foreground trước khi xếp hạng trên toàn ảnh.}
\end{table*}

\FloatBarrier
\subsubsection{Mở rộng theo công trình năm 2026 --- Prototype feedback với AlignG}
\textit{Nếu cho prototype của mỗi predicate thích nghi với ngữ cảnh của
từng ảnh, mean Recall có tăng so với cùng mô hình PE-Net hay không?}
Jung và Choi đề xuất cơ chế \textit{prototype feedback} trong AlignG
(ICML 2026)~\cite{jung2026aligng}. Nhóm tích hợp cơ chế này vào PE-Net dùng
đặc trưng CLIP đã cache. Trong mỗi ảnh, attention tổng hợp thông tin từ các
cặp đối tượng để GRU cập nhật prototype. Các prototype theo ảnh phản hồi
vào biểu diễn của từng cặp trước khi mô hình chấm điểm bằng cosine với
prototype toàn cục (Hình~\ref{fig:gqa-penet-aligng}).

\begin{table}[!ht]\centering\small
\caption{Ảnh hưởng của AlignG đến GC mR@50 trên 10.696 ảnh GQA test,
trung bình qua ba seed của Bảng~\ref{tab:gqatest}.}
\label{tab:gqa-aligng-test}
\begin{tabular}{@{}lrrr@{}}\toprule
Cấu hình nền & Tắt (\%) & Bật (\%) & $\Delta$ (điểm \%)\\\midrule
PE-CE   & 2,93 & 3,52 & $+0{,}59$\\
PE-PLPR & 4,07 & \textbf{4,62} & $+0{,}55$\\\bottomrule
\end{tabular}\\[2pt]
{\footnotesize GC = \textit{Graph Constraint}.}
\end{table}

Bật AlignG làm tăng mR@50 trung bình trên cả hai cấu hình nền
(Bảng~\ref{tab:gqa-aligng-test}). Trên nền PE-PLPR, hiệu mR@50 ghép theo
seed 42, 123 và 3407 lần lượt là $+1{,}05$, $+0{,}65$ và $-0{,}05$ điểm
phần trăm. Để kiểm tra riêng bước cập nhật prototype, cấu hình
\textit{Edge-only} giữ nhánh phản hồi đến cặp nhưng dùng prototype toàn
cục. Cấu hình này đạt 4,31\% mR@50. Khi bật cả bước cập nhật theo ảnh,
PE-PLPR + AlignG đạt $4{,}62 \pm 0{,}57$\% mR@50.

Khi thêm AlignG vào PE-PLPR, recall trung bình của hai lớp trái/phải gần
như giữ nguyên (86,59\% và 86,54\%), trong khi trung bình của 308 lớp còn
lại tăng từ 3,53 lên 4,09\%. Nhóm lớp này đóng góp phần tăng mR@50
quan sát được khi bật AlignG.

\subsubsection{Phân tích dự đoán trên validation}
Nhóm chọn PE-PLPR + AlignG để xem cấu hình dẫn đầu về mean Recall tìm
thấy và bỏ sót những quan hệ nào. Phụ lục~\ref{app:gqacases} trình bày
năm ca Hit và năm ca Miss theo GC top-50 với seed cố định 42,
kèm dự đoán của PE-PLPR cùng seed trên từng cặp.

Trên ảnh 1040, AlignG tìm được \textit{sign--on--window} ở hạng 7,
nhưng với \textit{balls--window} lại chọn \textit{to the left of} ở hạng 12
trong khi chú thích cần tìm là \textit{near}. Trên ảnh 1060, mô hình
chọn đúng \textit{window--to the right of--curtains} nhưng cặp này chỉ
xếp hạng 63. Ảnh 1037 minh họa tác dụng của GC rõ hơn vì cùng một cặp
\textit{pan--dough} có hai nhãn chuẩn \textit{to the left of} và \textit{of}.
Mô hình giữ nhãn thứ nhất ở hạng 40, nên chỉ một trong hai chú thích được
tính là Hit. Ở cả mười ví dụ, PE-PLPR và PE-PLPR + AlignG giữ cùng predicate. Khác biệt về thứ hạng chưa làm thay đổi kết quả Hit/Miss tại $K=50$.

\subsubsection{Kết luận Task A}
Trên dataset thứ hai, nhóm đã điều chỉnh đầu vào của PE-Net và đánh giá
theo cùng bài toán PredCls như trên VG150. PE-PL đạt tỷ lệ quan hệ chuẩn được tìm
thấy trong top-50 cao nhất khi lấy trung bình theo ảnh (R@50 = 78,21\%).
Khi lấy trung bình đều theo lớp, PE-PLPR + AlignG đứng đầu
(mR@50 = 4,62\%). Ngay cả với cấu hình này, recall trung bình của 308 lớp
ngoài hai quan hệ vị trí phổ biến nhất chỉ đạt 4,09\%.
Với ablation AlignG trên GQA, nhóm đề nghị xét điểm cộng theo tiêu chí
\textit{áp dụng công trình mới năm 2026} tại mục 5 của đề bài.
